\documentclass[12pt,letterpaper]{article}

% ------------------------------------------------
% Paquetes
% ------------------------------------------------
\usepackage[utf8]{inputenc}
\usepackage[T1]{fontenc}
\usepackage[spanish]{babel}

\usepackage{amsmath,amssymb,amsthm,mathtools,bm}
\usepackage{booktabs}
\usepackage{array}
\usepackage{geometry}
\usepackage{setspace}
\usepackage{microtype}
\usepackage{parskip}
\usepackage{titlesec}
\usepackage{fancyhdr}
\usepackage{enumitem}
\usepackage{xcolor}
\usepackage{listings}
\usepackage{tcolorbox}
\usepackage{hyperref}

\tcbuselibrary{skins,breakable}

% ------------------------------------------------
% Geometría
% ------------------------------------------------
\geometry{
  top=2.8cm,
  bottom=3cm,
  left=3.2cm,
  right=3.2cm,
  headheight=14pt
}

\setstretch{1.20}

% ------------------------------------------------
% Encabezados
% ------------------------------------------------
\pagestyle{fancy}
\fancyhf{}
\fancyhead[C]{\small\textsc{Econometría --- Gujarati --- Ejercicio 13.9}}
\fancyfoot[C]{\small\thepage}

\renewcommand{\headrulewidth}{0.4pt}

% ------------------------------------------------
% Títulos
% ------------------------------------------------
\titleformat{\section}
{\large\bfseries\scshape}
{\thesection.}{0.6em}{}
[\vspace{-4pt}\rule{\linewidth}{0.4pt}]

\titleformat{\subsection}
{\normalsize\bfseries}
{\thesubsection.}{0.5em}{}

% ------------------------------------------------
% Caja destacada
% ------------------------------------------------
\newtcolorbox{clave}{
  colback=white,
  colframe=black,
  boxrule=0.5pt,
  arc=3pt,
  left=8pt,
  right=8pt,
  top=6pt,
  bottom=6pt,
  breakable
}

% ------------------------------------------------
% Comandos útiles
% ------------------------------------------------
\newcommand{\Var}{\operatorname{Var}}
\newcommand{\Cov}{\operatorname{Cov}}
\newcommand{\E}{\operatorname{E}}
\newcommand{\plim}{\operatorname{plim}}

% ==========================================================
\begin{document}
% ==========================================================

\begin{center}

{\LARGE\bfseries Ejercicio 13.9 --- Error de Medición}\\[4pt]

{\LARGE\bfseries en el Modelo CAPM}\\[12pt]

{\large\itshape Econometría --- Gujarati \& Porter}\\[10pt]

\rule{0.55\linewidth}{0.6pt}\\[10pt]

{\normalsize
Ejercicio 13.9 resuelto por Emanuel Quintana Silva}\\[4pt]

{\small
Modelo de asignación de precios de activos de capital
y errores de medición en el coeficiente Beta}

\end{center}

\vspace{12pt}

% ==========================================================
\section{Enunciado}
% ==========================================================

El modelo de asignación de precios de activos de capital (CAPM) postula
la siguiente relación entre la tasa de rendimiento promedio de un valor
y su riesgo sistemático:

\begin{equation}
\bar{R}_{i} = \alpha_{1} + \alpha_{2}\beta_{i} + u_{i}
\end{equation}

donde:

\begin{itemize}
\item $\bar{R}_{i}$ = tasa de rendimiento promedio del valor $i$
\item $\beta_{i}$ = coeficiente Beta verdadero del valor $i$
\item $u_{i}$ = término de perturbación estocástico
\end{itemize}

La verdadera $\beta_{i}$ no es directamente observable. Se estima
mediante la \emph{línea característica}:

\begin{equation}
r_{it} = \alpha_{1} + \beta^{*} r_{mt} + e_{t}
\end{equation}

donde $r_{it}$ es la tasa de rendimiento del valor $i$ en el período
$t$ y $r_{mt}$ es la tasa de rendimiento del mercado (p.\,ej., índice
S\&P 500).

En la práctica se estima, por tanto:

\begin{equation}
\bar{R}_{i} = \alpha_{1} + \alpha_{2}\beta_{i}^{*} + u_{i}
\end{equation}

donde $\beta_{i}^{*}$ es el estimador de la Beta verdadera, y la
relación entre ambas es:

\begin{equation}
\beta_{i}^{*} = \beta_{i} + v_{i}
\end{equation}

siendo $v_{i}$ el \textbf{error de medición}.

\bigskip

\noindent Las preguntas del ejercicio son:

\begin{enumerate}[leftmargin=1.8em,label=\alph*)]
\item ¿Cuál será el efecto de este error de medición sobre la
      estimación de $\alpha_{2}$?
\item ¿El $\hat{\alpha}_{2}$ estimado de la ecuación (3) es una
      estimación insesgada del verdadero $\alpha_{2}$? De no ser así,
      ¿es consistente? De no ser así, ¿qué medidas correctivas sugiere?
\end{enumerate}

% ==========================================================
\section{Modelo observado con error de medición}
% ==========================================================

El modelo verdadero es la ecuación (1). Despejamos $\beta_{i}$ de (4):

\begin{equation}
\beta_{i} = \beta_{i}^{*} - v_{i}
\end{equation}

Sustituimos en el modelo verdadero:

\begin{equation}
\bar{R}_{i}
= \alpha_{1} + \alpha_{2}(\beta_{i}^{*} - v_{i}) + u_{i}
\end{equation}

Reordenando:

\begin{equation}
\bar{R}_{i}
= \alpha_{1} + \alpha_{2}\beta_{i}^{*}
  + \underbrace{(u_{i} - \alpha_{2} v_{i})}_{z_{i}}
\end{equation}

Definimos el \textbf{término de error compuesto}:

\begin{equation}
z_{i} = u_{i} - \alpha_{2} v_{i}
\end{equation}

de modo que el modelo estimado efectivamente es:

\begin{equation}
\bar{R}_{i} = \alpha_{1} + \alpha_{2}\beta_{i}^{*} + z_{i}
\end{equation}

Este modelo tiene la misma forma que la ecuación (3), pero el error
ya no es $u_{i}$ sino el compuesto $z_{i}$, que depende de $v_{i}$.

% ==========================================================
\section{Correlación entre regresor y error compuesto}
% ==========================================================

El supuesto fundamental de MCO exige:

\begin{equation}
\Cov(\beta_{i}^{*},\, z_{i}) = 0
\end{equation}

Calculemos esta covarianza. Como $\beta_{i}^{*} = \beta_{i} + v_{i}$
y $z_{i} = u_{i} - \alpha_{2} v_{i}$, y asumiendo que $u_{i}$,
$v_{i}$ y $\beta_{i}$ son mutuamente independientes:

\begin{align}
\Cov(\beta_{i}^{*},\, z_{i})
&= \Cov\!\left(\beta_{i} + v_{i},\; u_{i} - \alpha_{2} v_{i}\right)
\end{align}

Expandiendo:

\begin{align}
&= \Cov(\beta_{i}, u_{i})
  - \alpha_{2}\Cov(\beta_{i}, v_{i})
  + \Cov(v_{i}, u_{i})
  - \alpha_{2}\Cov(v_{i}, v_{i})
\end{align}

Usando los supuestos de independencia:

\begin{align}
\Cov(\beta_{i}, u_{i}) &= 0 \\
\Cov(\beta_{i}, v_{i}) &= 0 \\
\Cov(v_{i}, u_{i})     &= 0
\end{align}

queda únicamente:

\begin{equation}
\boxed{
\Cov(\beta_{i}^{*},\, z_{i})
= -\alpha_{2}\,\Var(v_{i})
= -\alpha_{2}\,\sigma_{v}^{2}
}
\end{equation}

Por lo tanto, siempre que $\sigma_{v}^{2} > 0$:

\begin{equation}
\Cov(\beta_{i}^{*},\, z_{i}) \neq 0
\end{equation}

El supuesto de exogeneidad del regresor queda violado y MCO
produce estimadores \textbf{sesgados e inconsistentes}.

% ==========================================================
\section{Inciso a) --- Efecto sobre la estimación de $\alpha_{2}$}
% ==========================================================

\subsection{Sesgo de atenuación}

La probabilidad límite del estimador MCO es:

\begin{equation}
\plim(\hat{\alpha}_{2})
= \alpha_{2}
  + \frac{\Cov(\beta_{i}^{*},\, z_{i})}{\Var(\beta_{i}^{*})}
\end{equation}

Ya obtuvimos:

\begin{equation}
\Cov(\beta_{i}^{*},\, z_{i}) = -\alpha_{2}\,\sigma_{v}^{2}
\end{equation}

Calculamos $\Var(\beta_{i}^{*})$. Como:

\begin{equation}
\beta_{i}^{*} = \beta_{i} + v_{i}
\end{equation}

y $\beta_{i}$ y $v_{i}$ son independientes:

\begin{align}
\Var(\beta_{i}^{*})
&= \Var(\beta_{i}) + \Var(v_{i}) \\[4pt]
&= \sigma_{\beta}^{2} + \sigma_{v}^{2}
\end{align}

Sustituyendo:

\begin{align}
\plim(\hat{\alpha}_{2})
&= \alpha_{2}
   + \frac{-\alpha_{2}\,\sigma_{v}^{2}}{\sigma_{\beta}^{2} + \sigma_{v}^{2}}
\end{align}

Factorizando:

\begin{align}
&= \alpha_{2}
   \left(
     1 - \frac{\sigma_{v}^{2}}{\sigma_{\beta}^{2} + \sigma_{v}^{2}}
   \right)
\end{align}

Simplificando:

\begin{equation}
\boxed{
\plim(\hat{\alpha}_{2})
= \alpha_{2}
  \left(
    \frac{\sigma_{\beta}^{2}}{\sigma_{\beta}^{2} + \sigma_{v}^{2}}
  \right)
}
\end{equation}

Como el factor entre paréntesis satisface:

\begin{equation}
0
<
\frac{\sigma_{\beta}^{2}}{\sigma_{\beta}^{2} + \sigma_{v}^{2}}
< 1
\end{equation}

se cumple:

\begin{equation}
\left|\plim(\hat{\alpha}_{2})\right| < |\alpha_{2}|
\end{equation}

\begin{clave}

El estimador MCO de $\alpha_{2}$ sufre de \textbf{sesgo de atenuación}
(\textit{attenuation bias}): el efecto estimado del riesgo sistemático
($\beta$) sobre el rendimiento esperado se \textbf{subestima
sistemáticamente} en valor absoluto, desplazándose artificialmente
hacia cero.

\end{clave}

\subsection{Dirección del sesgo}

Si la prima de riesgo es positiva ($\alpha_{2} > 0$), entonces:

\begin{equation}
\hat{\alpha}_{2} < \alpha_{2}
\end{equation}

Esto significa que la regresión \emph{subestima} cuánto aumenta el
rendimiento esperado por cada unidad adicional de riesgo sistemático.

% ==========================================================
\section{Inciso b) --- Insesgadez, consistencia y remedios}
% ==========================================================

\subsection{¿Es insesgado $\hat{\alpha}_{2}$?}

\textbf{No.}

La condición necesaria para la insesgadez es:

\begin{equation}
\E[z_{i} \mid \beta_{i}^{*}] = 0
\end{equation}

Pero hemos mostrado que:

\begin{equation}
\Cov(\beta_{i}^{*},\, z_{i}) = -\alpha_{2}\,\sigma_{v}^{2} \neq 0
\end{equation}

lo que implica que esta condición no se satisface. El error de medición
en la variable independiente destruye la insesgadez del estimador MCO
incluso en muestras pequeñas.

\subsection{¿Es consistente $\hat{\alpha}_{2}$?}

\textbf{No.}

A diferencia de lo que ocurre con errores en la variable dependiente
(que solo restan eficiencia sin afectar consistencia), el sesgo
provocado por errores de medición en el \emph{regresor} no desaparece
al aumentar el tamaño de la muestra. Formalmente:

\begin{equation}
\plim(\hat{\alpha}_{2})
= \alpha_{2}
  \underbrace{
    \left(
      \frac{\sigma_{\beta}^{2}}{\sigma_{\beta}^{2} + \sigma_{v}^{2}}
    \right)
  }_{<\,1}
  \neq \alpha_{2}
\end{equation}

siempre que $\sigma_{v}^{2} > 0$.

\begin{clave}

El estimador MCO es \textbf{inconsistente} en presencia de errores de
medición en la variable explicativa. Aumentar el tamaño de la muestra
no elimina el sesgo: el estimador converge a un valor
\emph{distinto} del parámetro verdadero.

\end{clave}

\subsection{Medidas correctivas}

\subsubsection*{1. Variables Instrumentales (VI) / MC2E}

Es el remedio principal. Se busca una variable instrumental $Z$ que
satisfaga dos condiciones:

\begin{enumerate}[leftmargin=1.8em]
\item \textbf{Relevancia:} alta correlación con la $\beta_{i}$
      verdadera:
      \[
        \Cov(Z,\, \beta_{i}) \neq 0
      \]
\item \textbf{Exogeneidad:} independencia de los errores $v_{i}$ y
      $u_{i}$:
      \[
        \Cov(Z,\, v_{i}) = 0
        \qquad\text{y}\qquad
        \Cov(Z,\, u_{i}) = 0
      \]
\end{enumerate}

Un instrumento común en el contexto del CAPM es la Beta estimada en
un período no solapado con el de la regresión de la Etapa II, o
variables de clasificación sectorial del activo.

\subsubsection*{2. Agrupamiento en carteras (\textit{portfolio grouping})}

En finanzas empírica es práctica estándar agrupar los activos en
\textbf{carteras ordenadas por Beta}. Al promediar rendimientos y
Betas dentro de cada cartera, los errores de medición individuales
$v_{i}$ tienden a cancelarse por la ley de los grandes números,
reduciendo la varianza del error compuesto y mitigando el sesgo de
atenuación.

\subsubsection*{3. Mejora de la precisión en la Etapa I}

Aumentar la frecuencia de los datos utilizados en la regresión (2)
--- por ejemplo, datos diarios en lugar de mensuales --- produce
estimaciones $\beta_{i}^{*}$ con menor varianza de error. Al reducir
$\sigma_{v}^{2}$, el factor de atenuación se aproxima a 1 y el sesgo
disminuye.

% ==========================================================
\section{Síntesis comparativa con el Ejercicio 13.8}
% ==========================================================

Los Ejercicios 13.8 y 13.9 ilustran el mismo fenómeno en dos contextos
distintos. La tabla siguiente resume los elementos clave:

\begin{center}
\renewcommand{\arraystretch}{1.4}
\begin{tabular}{lcc}
\toprule
\textbf{Elemento}        & \textbf{Ej.\ 13.8 (Friedman)} & \textbf{Ej.\ 13.9 (CAPM)} \\
\midrule
Variable con error       & $X_{i}$ (ingreso)   & $\beta_{i}^{*}$ (riesgo)   \\
Parámetro afectado       & $\hat{\beta}$       & $\hat{\alpha}_{2}$         \\
Error compuesto          & $u_{i}-\beta v_{i}$ & $u_{i}-\alpha_{2} v_{i}$   \\
$\plim$ del estimador    & $\dfrac{\beta}{1+\sigma_{v}^{2}/\sigma_{X}^{2}}$
                         & $\dfrac{\alpha_{2}\,\sigma_{\beta}^{2}}{\sigma_{\beta}^{2}+\sigma_{v}^{2}}$ \\
Dirección del sesgo      & Hacia cero          & Hacia cero                 \\
¿Consistente?            & No                  & No                         \\
\bottomrule
\end{tabular}
\end{center}

Nótese que ambas expresiones son algebraicamente equivalentes:

\begin{equation}
\frac{\beta}{1+\dfrac{\sigma_{v}^{2}}{\sigma_{X}^{2}}}
=
\beta \cdot \frac{\sigma_{X}^{2}}{\sigma_{X}^{2}+\sigma_{v}^{2}}
\end{equation}

la diferencia es solo de notación.

% ==========================================================
\section{Conclusión}
% ==========================================================

\begin{clave}

El uso de la Beta estimada $\beta_{i}^{*}$ como proxy de la Beta
verdadera $\beta_{i}$ introduce una correlación entre el regresor
y el término de error compuesto $z_{i} = u_{i} - \alpha_{2} v_{i}$:

\begin{equation*}
\Cov(\beta_{i}^{*},\, z_{i}) = -\alpha_{2}\,\sigma_{v}^{2} \neq 0
\end{equation*}

Como consecuencia:

\begin{enumerate}[leftmargin=1.8em]
\item El estimador MCO de $\alpha_{2}$ es \textbf{sesgado} hacia cero
      (\textit{attenuation bias}).
\item El estimador es \textbf{inconsistente}: el sesgo persiste
      incluso cuando $n \to \infty$.
\item El límite de probabilidad es:
\end{enumerate}

\begin{equation*}
\plim(\hat{\alpha}_{2})
= \alpha_{2}
  \left(
    \frac{\sigma_{\beta}^{2}}{\sigma_{\beta}^{2} + \sigma_{v}^{2}}
  \right)
< \alpha_{2}
\end{equation*}

Los remedios principales son: \textbf{Variables Instrumentales},
\textbf{agrupamiento en carteras} y \textbf{mejora de la precisión}
en la estimación de Beta de la Etapa I.

\end{clave}

\vspace{12pt}

\noindent\rule{\linewidth}{0.4pt}

\vspace{10pt}

\begin{center}
\small\textsc{Referencias}
\end{center}

\vspace{6pt}

\noindent
Gujarati, D.~N., \& Porter, D.~C. (2010).
\textit{Econometría} (5.a ed.).
McGraw-Hill.

\vspace{6pt}

\noindent
Wooldridge, J.~M. (2016).
\textit{Introductory Econometrics} (6.a ed.).
Cengage.

\vspace{6pt}

\noindent
Black, F., Jensen, M.~C., \& Scholes, M. (1972).
The capital asset pricing model: Some empirical tests.
En M.~C. Jensen (Ed.),
\textit{Studies in the Theory of Capital Markets}
(pp.\ 79--121). Praeger.

% ==========================================================
\end{document}
